\chapter{Fase 4 -- Balanço hidrostático}
\label{cap:fase4}

\section{Visão geral}

A Fase~4 é a de maior risco numérico de todo o trabalho: transforma os
campos de temperatura, pressão, umidade e vento já interpolados
verticalmente (Fase~3) em um conjunto de variáveis prognósticas
mutuamente consistentes com a métrica vertical e com o balanço
hidrostático de referência do \mpasa{} (\cref{sec:hidrostatico_teoria})
-- densidade seca $\rho_{zz}$, temperatura potencial úmida
$\theta_m$, pressão total $p$, água precipitável, e vento vertical
diagnóstico $w$. O algoritmo, reproduzido aqui passo a passo, foi
extraído literalmente do bloco \texttt{config\_met\_interp} da rotina
\texttt{init\_atm\_case\_gfs} de referência, preservando deliberadamente
duas fórmulas de $\rho_{zz}$ aparentemente redundantes que na verdade não
são -- um ponto de risco identificado explicitamente durante a
implementação e discutido na \cref{sec:duas_formulas}.

\section{Diagnóstico inicial de exner, $\theta$ e $\rho_{zz}$}

A partir da pressão e temperatura já interpoladas (Fase~3), calcula-se,
para cada célula $i$ e nível $k$, a função de Exner
$\Pi = (p/p_0)^{R/c_p}$, converte-se temperatura em temperatura
potencial, e obtém-se uma primeira estimativa \emph{diagnóstica} da
densidade seca:
\begin{equation}
  \Pi(k,i) = \left(\frac{p(k,i)}{p_0}\right)^{R/c_p},
  \qquad
  \theta(k,i) = T(k,i)\left(\frac{p_0}{p(k,i)}\right)^{R/c_p},
  \label{eq:exner_theta}
\end{equation}
\begin{equation}
  \rho_{zz}^{(0)}(k,i) = \frac{1}{1+q_v(k,i)}\cdot
  \frac{p(k,i)}{R\,\Pi(k,i)\,\theta(k,i)\left[1+(R_v/R - 1)\,q_v(k,i)\right]},
  \label{eq:rho_diagnostico}
\end{equation}
com $p_0=\SI{e5}{\pascal}$, $R=\SI{287}{\joule\per\kilo\gram\per\kelvin}$,
$c_p = \tfrac{7}{2}R$, e $R_v/R$ a razão entre as constantes de gás do
vapor d'água e do ar seco. Essa densidade diagnóstica é usada apenas para
a água precipitável (\cref{eq:precipitavel}) e para o acoplamento inicial
com o estado de referência -- \emph{não} é a densidade final do problema,
que exige resolver a relação hidrostática de forma consistente
(\cref{sec:iteracao}).

\begin{equation}
  W = \sum_{k} \rho_{zz}^{(0)}(k,i)\, q_v(k,i)\, \big[z_{grid}(k{+}1,i)-z_{grid}(k,i)\big].
  \label{eq:precipitavel}
\end{equation}

\section{Estado de referência isotérmico seco}

O estado base $(\,\overline{p},\overline{\Pi},\overline{\rho},\overline{\theta}\,)$
é uma atmosfera hipotética isotérmica seca a $T_0=\SI{250}{\kelvin}$,
avaliada no ponto médio geométrico de cada camada $z_m(k,i) =
\tfrac12[z_{grid}(k{+}1,i)+z_{grid}(k,i)]$:
\begin{align}
  \overline{p}(k,i) &= p_0\, \exp\!\left(-\frac{g\, z_m(k,i)}{R\,T_0}\right), \label{eq:estado_base_p}\\
  \overline{\Pi}(k,i) &= \left(\overline{p}(k,i)/p_0\right)^{R/c_p},
  \qquad
  \overline{\rho}(k,i) = \frac{\overline{p}(k,i)}{R\,T_0},
  \qquad
  \overline{\theta}(k,i) = \frac{T_0}{\overline{\Pi}(k,i)}. \label{eq:estado_base_resto}
\end{align}
Em seguida, tanto o estado base $\overline{\rho}$ quanto a densidade
diagnóstica $\rho_{zz}^{(0)}$ são \emph{acoplados à métrica vertical}
(divididos por $zz(k,i)$, \cref{eq:zgrid_zz}), e definem-se as
perturbações de referência
$p'(k,i) = p(k,i) - \overline{p}(k,i)$ e
$\rho'(k,i) = \rho_{zz}^{(0)}(k,i) - \overline{\rho}(k,i)$.

\section{Duas fórmulas de $\rho_{zz}$: por que não são redundantes}
\label{sec:duas_formulas}

\begin{caixaachado}
Um ponto de risco identificado explicitamente durante a implementação --
e preservado deliberadamente, sem tentativa de ``simplificação'' -- é
que o nível mais baixo ($k=1$) e os níveis subsequentes usam
\emph{duas fórmulas distintas} para $\rho_{zz}$, com expoentes e fatores
de correção de umidade diferentes:
\begin{equation}
  \rho_{zz}(1,i) = \frac{1}{zz(1,i)}\cdot\frac{p_0}{R}\cdot
  \frac{\left(p(1,i)/p_0\right)^{c_v/c_p}}{T(1,i)\left[1+1{,}61\,q_v(1,i)\right]},
  \label{eq:rho_nivel1}
\end{equation}
usando o expoente $c_v/c_p$ (não $R/c_p$ como em
\cref{eq:rho_diagnostico}) e o fator de correção de umidade $1{,}61\,q_v$
aplicado diretamente sobre $T$ (não a forma $(R_v/R-1)\,q_v$ do cálculo
diagnóstico inicial). A mesma segunda fórmula -- com o mesmo expoente
$c_v/c_p$ e o mesmo fator $1{,}61\,q_v$ -- reaparece dentro da iteração
hidrostática (\cref{eq:rho_iteracao}) para os níveis $k\ge 2$. Uma
tentativa de ``reconciliar'' essas duas expressões em uma única fórmula
canônica, por parecerem redundantes à primeira vista, introduziria um
erro sistemático sutil no campo de densidade -- por isso ambas foram
preservadas exatamente como aparecem no código-fonte de referência,
documentadas explicitamente no código Fortran resultante para que futuras
manutenções não as \emph{simplifiquem} por engano.
\end{caixaachado}

\section{Solução hidrostática iterativa}
\label{sec:iteracao}

Para os níveis $k=2,\dots,n_{VertLevels}$, resolve-se, célula a célula,
um sistema não linear pela relação hidrostática discreta do \mpasa{},
iterando até convergência (limiar $|\Delta p'| < \SI{0.0001}{\pascal}$,
máximo de 30 iterações):
\begin{align}
  p'(k,i) &\leftarrow p'(k{-}1,i) - \Big[f_{zm}(k)\rho'(k,i) + f_{zp}(k)\rho'(k{-}1,i)\Big]\,g\,\Delta z_u(k) \nonumber\\
  &\qquad - \Big[f_{zm}(k)\rho_{zz}(k,i)q_v(k,i) + f_{zp}(k)\rho_{zz}(k{-}1,i)q_v(k{-}1,i)\Big]\,g\,\Delta z_u(k), \label{eq:pp_iteracao}\\[0.4em]
  p(k,i) &= p'(k,i) + \overline{p}(k,i), \qquad
  \Pi(k,i) = \left(p(k,i)/p_0\right)^{R/c_p}, \label{eq:p_iteracao}\\[0.4em]
  \rho_{zz}(k,i) &= \frac{1}{zz(k,i)}\cdot\frac{p(k,i)}{R\,\Pi(k,i)\,T(k,i)\big[1+1{,}61\,q_v(k,i)\big]},
  \qquad
  \rho'(k,i) = \rho_{zz}(k,i)-\overline{\rho}(k,i). \label{eq:rho_iteracao}
\end{align}
Os coeficientes $f_{zm}(k)$, $f_{zp}(k)$ e $\Delta z_u(k)$ são exatamente
os já calculados na Fase~2 (\cref{cap:fase2}) -- é essa reutilização
direta que garante a consistência entre a métrica vertical e o campo de
massa resultante, mencionada na \cref{sec:hidrostatico_teoria}. Ao final
da convergência em todos os níveis, a temperatura potencial é
transformada de $\theta$ (seca) para $\theta_m$ (potencial úmida,
variável prognóstica real do modelo) por
$\theta_m(k,i) = \theta(k,i)\big[1+1{,}61\,q_v(k,i)\big]$, e a
perturbação $\rho'$ é desacoplada da métrica vertical multiplicando de
volta por $zz(k,i)$.

\section{Vento normal ponderado por densidade e vento vertical diagnóstico}

O produto $ru$ nas arestas -- a variável de fluxo de massa efetivamente
usada pelo núcleo dinâmico, não apenas o vento $u$ -- é
\begin{equation}
  ru(k,e) = u(k,e)\cdot \tfrac12\big[\rho_{zz}(k,c_1(e)) + \rho_{zz}(k,c_2(e))\big].
  \label{eq:ru}
\end{equation}
O vento vertical diagnóstico $w$ é obtido por um balanço de massa: a
divergência horizontal do fluxo $ru$ ao redor de cada célula, projetada
verticalmente pelos termos de métrica $zb$ (\cref{cap:fase2}) e, dado que
o \textit{namelist} real usa \texttt{config\_theta\_adv\_order=3},
acrescida do termo de correção de terceira ordem via $zb^3$ ponderado
pelo coeficiente \texttt{config\_coef\_3rd\_order}~$=0{,}25$:
\begin{equation}
  \rho w\big|_{k} \propto \sum_{e \in \partial i}
  \pm\Big[f_{zm}(k)\,zz(k,i)+f_{zp}(k)\,zz(k{-}1,i)\Big]
  \Big[zb(k,\cdot,e) + \operatorname{sgn}(ru)\cdot 0{,}25\, zb^3(k,\cdot,e)\Big]\, ru_{proj}(k,e),
  \label{eq:rw}
\end{equation}
com o sinal e o índice do segundo argumento de $zb$/$zb^3$ dependendo de
qual das duas células vizinhas da aresta $e$ é a célula $i$ em questão
(convenção idêntica à já usada no cálculo de $zx_u$, Fase~2). O vento
vertical final é $w(k,i) = \rho w|_k \big/ [f_{zp}(k)\rho_{zz}(k{-}1,i) +
f_{zm}(k)\rho_{zz}(k,i)]$, com $w(1,\cdot)=w(n_{VertLevels}+1,\cdot)=0$
por condição de contorno (nível do chão e topo rígido, respectivamente).

\section{Pressão de superfície diagnóstica e campos finais}

Um último campo diagnóstico, $\texttt{surface\_pressure}$ -- distinto do
$\texttt{psfc}$ já calculado na Fase~3 --, é obtido por extrapolação de
segunda ordem a partir dos dois níveis mais baixos:
\begin{equation}
  \texttt{surface\_pressure}(i) = \frac{0{,}5\,g}{rdz_w(1)}
  \Big[1{,}25\,\rho_{zz}(1,i)(1+q_v(1,i)) - 0{,}25\,\rho_{zz}(2,i)(1+q_v(2,i))\Big] + p'(1,i) + \overline{p}(1,i).
  \label{eq:surface_pressure}
\end{equation}
Finalmente, os campos prognósticos de saída são
$\rho(k,i) = \rho_{zz}(k,i)\cdot zz(k,i)$ (densidade física, desacoplada
da métrica) e $\theta(k,i) = \theta_m(k,i)/[1+1{,}61\,q_v(k,i)]$
(restaurando a temperatura potencial seca a partir da úmida, para o
campo \texttt{theta} de saída do \texttt{init.nc}).
